\documentclass[12pt]{article}
\usepackage[margin=1in]{geometry}
\usepackage{booktabs,graphicx,makecell,threeparttable}

\begin{document}

\providecommand{\StructuralRobustnessPath}{appendix/structural-robustness}

\clearpage
\subsection{Campaign-Day Evolution of the Reduced-Form Effects}
\label{app:campaign-day-dynamics}

Adults could seek treatment on twelve consecutive campaign days. We therefore
examine when the experimental effects emerge, without using an instrumental-variable
or conditional-hazard specification. The analysis uses the originally randomized
Close/Far assignment and the same 9,805-person sample and covariate adjustment as
the reduced-form endpoint analysis. Table~\ref{tab:campaign-original-assignment}
reproduces that endpoint: treatment rows are effects relative to Control, while
the Control row reports predicted take-up levels.

\begin{table}[htbp]
  \centering
  \small
  \caption{Original distance-assignment intention-to-treat estimates}
  \label{tab:campaign-original-assignment}
  \input{\StructuralRobustnessPath/tables/rf-original-distance-itt.tex}
\end{table}

For campaign day $t$, cumulative take-up is an indicator for directly recorded
treatment on or before day $t$. Daily incidence is an indicator for directly
recorded treatment exactly on day $t$. Both outcomes retain the complete randomized
sample, so daily incidence is an unconditional probability rather than a hazard
among the post-treatment-selected set of people who have not yet participated.
Day 12 consequently reproduces the original endpoint exactly.

We estimate the four treatment arms interacted with original Close/Far assignment,
adjusting for gender, age, expected distance, and county fixed effects. Inference
uses 1,000 paired exponential cluster-weight draws at the 144 randomization clusters.
The same cluster weights are used for every campaign day and derived contrast.
Figure~\ref{fig:campaign-cumulative} reports adjusted cumulative take-up; ribbons
are pointwise 95 percent intervals.

\begin{figure}[htbp]
  \centering
  \includegraphics[width=\linewidth]{\StructuralRobustnessPath/figures/campaign-day-cumulative.pdf}
  \caption{Adjusted cumulative take-up over the campaign}
  \label{fig:campaign-cumulative}
  \begin{minipage}{0.94\linewidth}
    \vspace{0.3em}\footnotesize
    \textit{Notes:} Shaded regions are pointwise 95 percent intervals from
    paired exponential cluster-weight draws. All estimands retain the full
    randomized sample.
  \end{minipage}
\end{figure}

Table~\ref{tab:campaign-interactions} shows the treatment-by-Far interactions
at four campaign milestones. The pooled Any Signal--No Signal Far--Close effect
grows from 0.029 on day 1 to 0.082 on day 6 and 0.124 on day 12. Thus approximately
two-thirds of the endpoint differential has emerged by day 6. The arm-specific
results show that the pattern is clearest for Ink. The Bracelet interaction moves
in the same direction later in the campaign but remains imprecisely estimated,
while Calendar shows no comparable pattern.

\begin{table}[htbp]
  \centering
  \small
  \caption{Evolution of treatment-by-Far interaction effects}
  \label{tab:campaign-interactions}
  \resizebox{\linewidth}{!}{%
    \begin{tabular}{lcccc}
\toprule
 & Day 1 & Day 3 & Day 6 & Day 12 \\
\midrule
\multicolumn{5}{l}{\textit{Panel A: Arm-specific effects}} \\
Calendar $\times$ Far & \makecell[c]{0.003\\{[-0.039, 0.046]}} & \makecell[c]{-0.037\\{[-0.111, 0.035]}} & \makecell[c]{0.003\\{[-0.077, 0.083]}} & \makecell[c]{-0.010\\{[-0.110, 0.081]}} \\
Ink $\times$ Far & \makecell[c]{0.047\\{[0.004, 0.087]}} & \makecell[c]{0.071\\{[-0.002, 0.140]}} & \makecell[c]{0.143\\{[0.049, 0.237]}} & \makecell[c]{0.172\\{[0.048, 0.289]}} \\
Bracelet $\times$ Far & \makecell[c]{0.015\\{[-0.028, 0.059]}} & \makecell[c]{-0.010\\{[-0.080, 0.057]}} & \makecell[c]{0.024\\{[-0.056, 0.111]}} & \makecell[c]{0.065\\{[-0.042, 0.177]}} \\
\addlinespace
\multicolumn{5}{l}{\textit{Panel B: Pooled signal effect}} \\
Any Signal $\times$ Far & \makecell[c]{0.029\\{[0.001, 0.059]}} & \makecell[c]{0.049\\{[-0.002, 0.099]}} & \makecell[c]{0.082\\{[0.023, 0.143]}} & \makecell[c]{0.124\\{[0.052, 0.199]}} \\
\bottomrule
\end{tabular}
}
  \begin{minipage}{0.94\linewidth}
    \vspace{0.3em}\footnotesize
    \textit{Notes:} Each arm-specific row is the treatment-versus-Control effect
    under Far assignment minus the corresponding effect under Close assignment.
    The pooled row compares the average of Ink and Bracelet with the average of
    Control and Calendar. Brackets contain pointwise 95 percent cluster-weight
    intervals.
  \end{minipage}
\end{table}

For completeness, Table~\ref{tab:campaign-cumulative-levels} reports adjusted
levels and within-arm Far-minus-Close differences on the same milestone days.

\begin{table}[htbp]
  \centering
  \small
  \caption{Adjusted cumulative take-up on selected campaign days}
  \label{tab:campaign-cumulative-levels}
  \resizebox{\linewidth}{!}{%
    \begin{tabular}{lcccc}
\toprule
 & Day 1 & Day 3 & Day 6 & Day 12 \\
\midrule
\multicolumn{5}{l}{\textit{Panel A: Close assignment}} \\
Control & \makecell[c]{0.111\\{[0.083, 0.139]}} & \makecell[c]{0.207\\{[0.167, 0.244]}} & \makecell[c]{0.327\\{[0.280, 0.371]}} & \makecell[c]{0.440\\{[0.389, 0.487]}} \\
Calendar & \makecell[c]{0.099\\{[0.072, 0.129]}} & \makecell[c]{0.238\\{[0.194, 0.284]}} & \makecell[c]{0.338\\{[0.289, 0.388]}} & \makecell[c]{0.456\\{[0.423, 0.493]}} \\
Ink & \makecell[c]{0.067\\{[0.044, 0.096]}} & \makecell[c]{0.145\\{[0.108, 0.187]}} & \makecell[c]{0.200\\{[0.151, 0.250]}} & \makecell[c]{0.313\\{[0.259, 0.370]}} \\
Bracelet & \makecell[c]{0.102\\{[0.073, 0.134]}} & \makecell[c]{0.242\\{[0.201, 0.291]}} & \makecell[c]{0.361\\{[0.306, 0.416]}} & \makecell[c]{0.474\\{[0.406, 0.535]}} \\
\addlinespace
\multicolumn{5}{l}{\textit{Panel B: Far assignment}} \\
Control & \makecell[c]{0.035\\{[0.024, 0.047]}} & \makecell[c]{0.090\\{[0.063, 0.119]}} & \makecell[c]{0.153\\{[0.115, 0.190]}} & \makecell[c]{0.259\\{[0.204, 0.323]}} \\
Calendar & \makecell[c]{0.027\\{[0.018, 0.036]}} & \makecell[c]{0.085\\{[0.062, 0.107]}} & \makecell[c]{0.166\\{[0.142, 0.190]}} & \makecell[c]{0.265\\{[0.232, 0.297]}} \\
Ink & \makecell[c]{0.038\\{[0.025, 0.051]}} & \makecell[c]{0.099\\{[0.068, 0.134]}} & \makecell[c]{0.169\\{[0.116, 0.226]}} & \makecell[c]{0.305\\{[0.240, 0.362]}} \\
Bracelet & \makecell[c]{0.041\\{[0.028, 0.055]}} & \makecell[c]{0.116\\{[0.091, 0.145]}} & \makecell[c]{0.210\\{[0.176, 0.242]}} & \makecell[c]{0.358\\{[0.314, 0.408]}} \\
\addlinespace
\multicolumn{5}{l}{\textit{Panel C: Far minus Close}} \\
Control & \makecell[c]{-0.076\\{[-0.105, -0.049]}} & \makecell[c]{-0.117\\{[-0.163, -0.071]}} & \makecell[c]{-0.175\\{[-0.230, -0.121]}} & \makecell[c]{-0.181\\{[-0.254, -0.098]}} \\
Calendar & \makecell[c]{-0.073\\{[-0.104, -0.043]}} & \makecell[c]{-0.154\\{[-0.204, -0.102]}} & \makecell[c]{-0.172\\{[-0.227, -0.116]}} & \makecell[c]{-0.191\\{[-0.241, -0.142]}} \\
Ink & \makecell[c]{-0.029\\{[-0.060, -0.002]}} & \makecell[c]{-0.046\\{[-0.102, 0.006]}} & \makecell[c]{-0.031\\{[-0.106, 0.044]}} & \makecell[c]{-0.008\\{[-0.091, 0.076]}} \\
Bracelet & \makecell[c]{-0.061\\{[-0.096, -0.028]}} & \makecell[c]{-0.126\\{[-0.179, -0.075]}} & \makecell[c]{-0.151\\{[-0.214, -0.088]}} & \makecell[c]{-0.116\\{[-0.192, -0.029]}} \\
\bottomrule
\end{tabular}
}
  \begin{minipage}{0.94\linewidth}
    \vspace{0.3em}\footnotesize
    \textit{Notes:} Cells report adjusted predictions or Far-minus-Close
    differences. Brackets contain pointwise 95 percent cluster-weight intervals.
  \end{minipage}
\end{table}

Figure~\ref{fig:campaign-incidence} shows the corresponding unconditional daily
take-up probabilities. The differential effect of public signals under Far
assignment develops over the campaign rather than appearing only as an endpoint
contrast. This accumulation is consistent with the value of visible signals
increasing as participation unfolds within communities and provides additional
support for the broad social-image or social-information mechanism. Timing alone
does not, however, separately identify social image from related mechanisms such
as reminders, descriptive norms, or social learning.

\begin{figure}[htbp]
  \centering
  \includegraphics[width=\linewidth]{\StructuralRobustnessPath/figures/campaign-day-incidence.pdf}
  \caption{Adjusted unconditional daily take-up}
  \label{fig:campaign-incidence}
  \begin{minipage}{0.94\linewidth}
    \vspace{0.3em}\footnotesize
    \textit{Notes:} Shaded regions are pointwise 95 percent cluster-weight
    intervals. These are unconditional daily probabilities, not hazards among
    those who have not yet participated.
  \end{minipage}
\end{figure}

\clearpage
\end{document}
